\section{Related Work and Prior Transfer Evidence}
\label{sec:related}

\paragraph{Structural difficulty and prior Tool-to-KG evidence.}
Tool Graph Complexity distinguishes reference-set non-minimality from dependency composition and analyzes motif-associated tool-selection difficulty~\cite{anonymousToolGraph}. Its adjusted analysis emphasizes sequential-pattern penalties, whereas raw motif rankings also reflect size and surface properties~\cite{anonymousToolGraph}. Appendix~O already evaluates entity selection on a finance KG, including motif and size variation, explicit relations, and adjacent decoys; it reports easier chains with explicit KG relations, difficult incoming trees, and decoy-related inclusion errors~\cite{anonymousToolGraph}. These are antecedent findings for our diagnostic follow-up, which separates relation disclosure from matched-fact presentation and examines residual selection errors. Differences in references, prompts, populations, and evaluation protocols limit direct score comparisons; unquantified task/response overlap also precludes a claim of independent replication.

\paragraph{Graph structure and task comprehension.}
Talk like a Graph evaluates how graph encodings, task types, and graph structure affect LLM performance~\cite{fatemi2024talklikegraph}. Xinnan Dai et al. study pattern comprehension and discovery from terminological and topological descriptions, including pattern detection and graph modification~\cite{dai2025graphpatterns}. GraphOmni jointly varies graph families, serialization formats, and prompting schemes across graph-theoretic tasks~\cite{xu2026graphomni}. Our motifs organize comparisons of query-conditioned entity selection: the reference includes intermediate and branching entities on all permitted supporting routes. The analysis concerns recovery of this declared route reference, with structure interpreted alongside query semantics and visible evidence.

\paragraph{Availability and representation of relational evidence.}
Xinbang Dai et al. compare KG input formats and triple organization and evaluate robustness to noisy or incomplete subgraphs using answer-based evaluation~\cite{dai2024kgunderstanding}. Tian et al. distinguish graph transformation from linearization and examine formats, orders, and templates in KG-augmented generation~\cite{tian2025alignment}. We distinguish adding explicit relations to native entity cards from changing the presentation of designated matched facts. Raw versus Raw+E changes available information; controlled triples versus templated prose and source-, destination-, or both-card hosting change how that information is presented. Native descriptions are not assumed to express every graph relation, so disclosure gains alone cannot establish better reasoning from equally informative inputs.

\paragraph{Graph retrieval and route recovery.}
STaRK evaluates retrieval of answer entities satisfying textual and relational query constraints~\cite{wu2024stark}. G-Retriever constructs graph context using prize-collecting Steiner tree optimization, while Reasoning on Graphs plans relation sequences and retrieves instantiated reasoning paths for KG question answering~\cite{he2024gretriever,luo2024rog}. GraphWalk studies tool-based graph navigation and reports incomplete exploration in its graph-in-context setting~\cite{ghandi2026graphwalk}. Our main evaluation holds target availability within a supplied candidate pool fixed and scores selected entity sets against the complete route reference. This objective differs from retrieving answer entities, constructing a compact useful subgraph, or generating an answer from retrieved paths~\cite{wu2024stark,he2024gretriever,luo2024rog}. Final-set omissions are observable selection errors; they do not reconstruct internal traversal steps or identify a shared failure mechanism across these settings.

\paragraph{Distractors and selection metrics.}
Amiraz et al. quantify how irrelevant passages can distract answer-generating models, distinguishing distraction from a binary relevance label~\cite{amiraz2025distracting}. TopoRAG addresses noise from heuristic neighbor expansion through retrieval constrained by induced-subgraph diameter~\cite{wu2026toporag}. Our disconnected/adjacent-context comparison instead examines query-irrelevant neighboring entities within a fixed-size pool containing every reference entity. Structural proximity is a different distractor construction from passage-level semantic distraction and does not by itself establish irrelevance~\cite{amiraz2025distracting}. Li and Ouyang show that the relationship between knowledge-selection quality and generation quality depends on the generator and task~\cite{li2025selection}. Accordingly, we use recall to measure route coverage and precision and exact recovery to diagnose excess selection and complete-set recovery; these metrics establish neither minimum sufficient evidence nor downstream answer correctness.
